% =====================================================================
% AutoRec with a heteroscedastic uncertainty head (week4/autoRec2.ipynb)
% All numbers below are from the notebook's own saved run on the
% CORRECTED ML-100K (943 users, 1,682 items, 70k/20k/10k split) -
% not the ml-latest-small shape (610 users, 9,724 items) an earlier
% run of this same notebook used before the dataset-identity fix.
% =====================================================================

\subsection{AutoRec with a Learned Uncertainty Head}
\label{sec:autorec-uncertainty}

As a complementary probe of where predictive uncertainty comes from, we ran a
separate experiment on item-based AutoRec \citep{sedhain2015autorec}: does
uncertainty have to be inferred after the fact from a plain model's mistakes,
or is it worth training a second head to predict it directly? We compare
three ways of scoring how much to trust a prediction, plus the two
uncertainty-free abstention rules from Section~\ref{sec:abstention-rules}
(Algorithms 2 and 3), all evaluated as selective-prediction curves against a
random-abstention baseline and an oracle upper bound.

\subsubsection{Models}

\textbf{Vanilla AutoRec} is the standard item-based autoencoder: a sigmoid
encoder, dropout, and a linear decoder, trained with masked MSE over
observed ratings only. \textbf{Heteroscedastic AutoRec} shares the same
encoder but splits the decoder into two heads: a mean head $\mu(x)$ and a
variance head, $\mathrm{var}(x) = \mathrm{softplus}(\text{raw\_var}(x)) +
10^{-4}$, trained jointly via Gaussian negative log-likelihood,
\[
  \mathcal{L} = \tfrac{1}{2}\Big[\log \mathrm{var}(x) + \frac{(y-\mu(x))^2}{\mathrm{var}(x)}\Big],
\]
with gradient-norm clipping (max norm 5.0) to keep the variance head stable
early in training.

\textbf{Uncertainty signals compared:}
\begin{itemize}
  \item \emph{Residual (item)}: the vanilla model's mean squared
        reconstruction error on item $i$'s \emph{observed} train+val
        ratings, reused as a proxy score for every test rating of that
        item. Since I-AutoRec is item-indexed, this is free to compute but
        collapses every test rating of the same item to one identical
        score - it cannot distinguish which user is harder to predict.
  \item \emph{Residual (user+item)}: the same idea, but averaging the
        item-level score with an analogous user-level score,
        $0.5(\text{item\_score}[i] + \text{user\_score}[u])$, to make the
        signal target-specific.
  \item \emph{Heteroscedastic variance}: the trained variance head's output
        $\mathrm{var}(x)_{ui}$ directly - the only one of the three that is
        genuinely per-(user,item), rather than inherited from a group
        average.
\end{itemize}

\subsubsection{Hyperparameter search}

The backbone was tuned in three passes, all reusing the same 70/20/10
proposal split: (1) a coarse grid over
hidden~$\in\{50,100,200,300,500,750,1000,1500\}$~$\times$
lr~$\in\{5{\times}10^{-4},\ldots,1.5{\times}10^{-2}\}$~$\times$
weight\_decay~$\in\{0,10^{-5},10^{-4},10^{-3},3{\times}10^{-3},10^{-2}\}$
(384 configs); (2) a focused refinement fixing hidden$=500$ (the coarse
winner) and sweeping a denser local lr~$\times$~weight\_decay grid around
it; (3) a dropout sweep at the winning backbone. Table~\ref{tab:autorec2-hp}
shows each stage only ever improved on, never regressed, the previous
stage's winner - notably the focused search moved lr from $2{\times}10^{-3}$
to $2.5{\times}10^{-3}$ and improved validation RMSE further, confirming
hidden$=500$ is a genuine local optimum rather than an artifact of a coarse
grid's edge.

\begin{table}[htbp]
\centering
\caption{AutoRec backbone hyperparameter search, by stage}
\label{tab:autorec2-hp}
\begin{tabular}{lcccc}
\toprule
Stage & hidden & lr & weight\_decay & Val.\ RMSE \\
\midrule
1. Coarse grid (384 configs)      & 500 & 0.0020 & $1{\times}10^{-3}$ & 0.8955 \\
2. Focused search at hidden=500 (42 configs) & 500 & 0.0025 & $1{\times}10^{-3}$ & 0.8953 \\
3. Dropout sweep (10 values)      & 500 & 0.0025 & $1{\times}10^{-3}$ & 0.8946 \\
\bottomrule
\end{tabular}
\vspace{2pt}
\\ \footnotesize Final chosen config: hidden=500, lr=0.0025, weight\_decay=$1{\times}10^{-3}$, dropout=0.2.
\end{table}

\subsubsection{Headline results}

\begin{table}[htbp]
\centering
\caption{Test RMSE and uncertainty quality, refit on train+val, corrected ML-100K}
\label{tab:autorec2-headline}
\begin{tabular}{lccc}
\toprule
Model & Test RMSE & Calibration $\mathbb{E}[(y-\hat\mu)^2/\mathrm{var}]$ & Spearman $\rho$ (uncertainty, sq.\ error) \\
\midrule
Vanilla AutoRec         & 0.8855 & --- & --- \\
\quad residual (item)         & --- & --- & 0.109 \\
\quad residual (user+item)    & --- & --- & 0.202 \\
Heteroscedastic AutoRec & 0.8871 & 1.973 (ideal $\approx 1$) & 0.238 \\
\bottomrule
\end{tabular}
\end{table}

The heteroscedastic head's variance is the single best-correlated signal
with true error, but the calibration check shows it is meaningfully
\emph{overconfident} on the real (sparser, 943$\times$1{,}682) benchmark -
true squared error runs on average $1.97\times$ the predicted variance,
versus the near-ideal $1.05\times$ this same notebook found before the
dataset-identity fix on the denser, easier ml-latest-small shape. This is
the more informative of the two runs precisely because it is the harder
dataset: it shows the learned uncertainty head is directionally correct but
has real, fixable headroom (e.g.\ a warm-up schedule or a better variance
prior) rather than being already near its ceiling.

\subsubsection{Selective prediction}

\begin{table}[htbp]
\centering
\caption{Selective RMSE vs.\ coverage, all six abstention rules plus oracle}
\label{tab:autorec2-selective}
\resizebox{\textwidth}{!}{%
\begin{tabular}{lcccccc}
\toprule
Method & \multicolumn{6}{c}{Coverage (fraction of test ratings kept)} \\
\cmidrule(lr){2-7}
 & 1.00 & 0.90 & 0.80 & 0.70 & 0.60 & 0.50 \\
\midrule
Random abstention (baseline)         & 0.8855 & 0.8854 & 0.8860 & 0.8851 & 0.8850 & 0.8856 \\
Residual (item)                      & 0.8855 & 0.8671 & 0.8528 & 0.8408 & 0.8310 & 0.8271 \\
Low user support (Alg.\ 3)           & 0.8855 & 0.8772 & 0.8701 & 0.8676 & 0.8652 & 0.8606 \\
Low predicted rating (Alg.\ 2)       & 0.8855 & 0.8794 & 0.8588 & 0.8446 & 0.8285 & 0.8082 \\
Residual (user+item)                 & 0.8855 & 0.8512 & 0.8258 & 0.8045 & 0.7819 & 0.7631 \\
\textbf{Heteroscedastic variance}    & \textbf{0.8871} & \textbf{0.8452} & \textbf{0.8253} & \textbf{0.8036} & \textbf{0.7799} & \textbf{0.7596} \\
\midrule
Oracle (true-error ranking, unattainable) & 0.8855 & 0.6785 & 0.5631 & 0.4721 & 0.3909 & 0.3167 \\
\bottomrule
\end{tabular}%
}
\end{table}

\subsubsection{Discussion}

Two findings stand out. First, the ranking of methods is consistent and
sensible: both naive, uncertainty-free rules (Algorithms 2 and 3) beat
random abstention, but neither beats a method that actually looks at
reconstruction error - and \emph{residual (item)} being the weakest
non-random score confirms the expected failure mode, since it assigns every
test rating of a given item the exact same score regardless of who the user
is. Second, and more striking, the learned heteroscedastic head's edge over
the free \emph{residual (user+item)} heuristic - which is nothing more than
an average of already-computed training errors, requiring no extra
parameters or a second loss term - is thin: a 0.15\% RMSE improvement at 50\%
coverage (0.7596 vs.\ 0.7631) and a modest Spearman gain (0.238 vs.\ 0.202).
Combined with the calibration gap above, this suggests the learned signal is
not yet extracting much more information than is already implicit in the
model's own training residuals on this dataset - closing the calibration gap
is a concrete, testable next step for whether a trained uncertainty head can
earn a larger margin over the heuristic it is currently barely beating.
